\documentclass[11pt,a4paper]{article}
\usepackage[T1]{fontenc}
\usepackage[utf8]{inputenc}
\usepackage{lmodern}
\usepackage[margin=28mm]{geometry}
\usepackage{amsmath,amssymb,amsthm,mathtools}
\usepackage{enumitem}
\usepackage{booktabs,array}
\usepackage{xurl}
\usepackage[colorlinks=true,linkcolor=blue,urlcolor=blue]{hyperref}
\usepackage{microtype}
\setlength{\emergencystretch}{2em}
\newtheorem{theorem}{Theorem}
\newtheorem{lemma}[theorem]{Lemma}
\theoremstyle{remark}
\newtheorem*{remark}{Remark}
\newcommand{\N}{\mathbb N}
\newcommand{\Z}{\mathbb Z}
\newcommand{\Co}{C_{\mathrm{outer}}}
\newcommand{\Cs}{C_{\mathrm{scalar}}}
\newcommand{\Up}{\mathsf{Up}}
\newcommand{\Down}{\mathsf{Down}}
\newcommand{\tot}[1]{\operatorname{tot}(#1)}
\newcommand{\PA}{\mathrm{PA}}
\DeclareMathOperator{\proj}{proj}
\title{Decidability of Symmetric Three-Counter Reachability\\
under Assumptions S2 and S3}
\author{}
\date{1 November 2026\\Revised explanatory version}
\begin{document}
\maketitle

\begin{abstract}
Assuming effective semilinearity of reachability for symmetric
 two-counter machines and symmetric three-dimensional VASS, the argument
reduces three-counter reachability to terminating upward saturations.
The geometric part distinguishes a boundary face, which has two free
coordinates, from an axis collar, which has only one. A potential based
on the two smallest coordinates bounds the number of large excursions
when an endpoint is near an axis. Scalar pumping handles repeated axis
visits. The resulting algorithm needs only an emptiness test between
boxes for one fixed connector relation; it does not compute the full
three-counter reachability relation.
\end{abstract}

\paragraph{Scope of this version.}
This document integrates the explanation of the proof with the earlier
assembled argument. It spells out the roles of the interfaces, the pair
potential, the connector oracle, and the two saturation procedures.
The longer geometric and pumping constructions are summarized here and
referenced in the accompanying research log at the end. Thus this is an
explanatory research proof, not a machine-checked formalization or a
replacement for every supporting construction. No effective
semilinearity theorem for full three-counter reachability is asserted.

\tableofcontents
\bigskip
\section{The problem and the proof roadmap}\label{sec:roadmap}

A configuration is a control state together with a vector in \(\N^3\).
An update transition adds an integer vector and is enabled when the
resulting counters remain nonnegative. A zero-test transition checks
that its named counter is zero. A VASS has updates but no zero tests.

\begin{theorem}\label{thm:main}
Assume the following two algorithms.
\begin{description}[leftmargin=2.8em,style=nextline]
\item[\textnormal{\textbf{S2}}] Given a symmetric two-counter machine,
compute a Presburger formula for its full reachability relation.
\item[\textnormal{\textbf{S3}}] Given a symmetric three-dimensional VASS,
compute a Presburger formula for its full reachability relation.
\end{description}
Then point-to-point reachability in a symmetric three-counter machine
is decidable.
\end{theorem}


Here symmetry means closure of every update transition under permutations of its counter coordinates, preserving its source and target control states. Zero tests name their tested counter. The construction does not assume that transitions are reversible. Reversal is used to construct and analyze a different machine, which remains symmetric.

The difficulty is not computing the relation of one fixed segment:
\textbf{S2} and \textbf{S3} do that. The difficulty is taking unboundedly
many segments. The transitive closure of a semilinear relation need not
be semilinear, so we cannot simply compose all boundary relations and
invoke Presburger arithmetic.

The key geometric observation is that a zero-test face still has two
free coordinates, whereas a visit near an axis has only one. For example,
\[
  x=(0,u,v)
  \qquad\Longrightarrow\qquad
  s_2(x)=\min(u,v),
\]
where \(s_2\) is the sum of the two smallest coordinates. Thus a bound
on \(s_2\) measures proximity to an axis even when the total
\(u+v\) is arbitrarily large. This is the quantity used to control
excursions; a bound on total alone would not give the required argument.

Here is the logical order of the proof.
\begin{enumerate}[leftmargin=*]
\item \textbf{Describe the boundary interfaces.}
Store bounded coordinates in finite control. Faces become two-counter
charts; axis collars become one-counter charts
(Section~\ref{sec:ingredients}).
\item \textbf{Bound large excursions between axis visits.}
Expose actual independent pumps and transfer loops. In what remains,
a pair potential increases at every large positive excursion. A bounded
target pair therefore bounds their number. Reverse the argument for
negative excursions (Section~\ref{sec:pair}). Sign changes are normalized
using peaks, valleys, and carefully placed hub events
(Section~\ref{sec:normal}).
\item \textbf{Decide a fixed connector relation \(K\) between boxes.}
A two-dimensional target box is either upward or has a bounded
coordinate. The first case uses upward saturation; the second uses the
excursion bound. Repeated axis visits are handled by a separate scalar
saturation (Section~\ref{sec:oracle}).
\item \textbf{Add the exposed pumps back.}
Compute negative closure backwards and positive closure forwards.
Dickson's lemma makes both saturations terminate
(Section~\ref{sec:completion}).
\end{enumerate}
This route does not require semilinearity of the closure of all
bounded-gain boundary segments.

\paragraph{Notation.}
For a vector \(z\), write \(\tot{z}=\sum_i z_i\), including when its
entries are signed. All finite control labels are implicit in set and
relation notation. A semilinear set and a Presburger-definable set are
used interchangeably, with effective conversions.
We write PA for Presburger-definable. Relations are composed from left to
right: \(x(R;S)z\) means that \(xRy\) and \(ySz\) for some \(y\).
The symbol \(R^*\) denotes reflexive-transitive closure. For a set \(S\),
\(S;R\) is its forward image, and \(R^{-1}(T)\) is the preimage of \(T\).

\section{Interfaces and geometric ingredients}\label{sec:ingredients}

\subsection{Faces, axis collars, and the lower-dimensional machines}

Fix a machine-dependent bound \(H\). In a region where some physical
counter is at most \(H\), choose one such counter as a \emph{marker}.
Its index \(i\) and value \(b\) are stored in finite control. A face
mode therefore has the form \((q,i,b)\), and its two counters are the
other two physical values. These two values are not replaced by their
sum. This is the basic two-dimensional chart.

Later choose a larger width \(W\). An \emph{axis collar} consists of
configurations with two bounded physical coordinates. Store their
indices and values in control and keep the third value as a single
counter. The axis need not be literal: a point \((b,c,n)\), with fixed
bounded \(b,c\), is just as scalar as \((0,0,n)\). Exact fixed-value
adapters can move between these presentations when needed.

The proof uses the following different interfaces. Keeping their
meanings separate is essential.
\begin{center}
\small
\begin{tabular}{@{}p{.24\linewidth}cp{.58\linewidth}@{}}
\toprule
Interface & Dimension & Meaning of its coordinates\\
\midrule
Physical face & 2 & Two free physical counters; one marker is stored.\\
Physical axis collar & 1 & One free physical counter; two bounded values are stored.\\
High pump interface & 3 & All three physical values, above a fixed enabling threshold.\\
Transfer hub & 1 & Total and finite residue data at a selected transfer anchor.\\
Cartesian reset port & 0 & A private finite-control interface between an entrance and an exit relation.\\
\bottomrule
\end{tabular}
\end{center}
A scalar hub is therefore not automatically a physical axis. That
qualification will matter when applying the pair bound.

If an update has absolute coordinate size at most \(J\), changing the
marker from \(i\) to \(j\) forces both coordinates to be bounded at the
change: the old marker is at most \(H+J\) afterwards, and the new marker
is at most \(H+J\) beforehand. Such changes are axis-collar events once
\(W\ge H+J\). Between axis-collar events, the marker is fixed. A zero
test of a free coordinate also visits an axis collar. Consequently the
intervening fixed-marker corridor has no free-coordinate zero test and
is represented by a symmetric two-VASS. The whole thin subsystem,
including marker changes and zero tests, has a symmetric two-CM
presentation and is covered by \textbf{S2}.

\subsection{What the geometric constructions provide}

A \emph{macro relation} records the exact endpoint relation of a specified
family of actual runs. It is not a new arbitrary transition. Below,
\emph{positive} and \emph{negative} refer to the signs of cycle totals
in an update component. A \emph{large} excursion, also called a
\emph{fat edge}, is a selected macro whose endpoint total gain exceeds
the fixed cutoff \(\Co\) in a positive phase, or whose total loss
exceeds it in a negative phase. These two distinctions are different:
a positive component can also contain bounded-gain paths.

For an update component, its \emph{cycle cone} is the nonnegative real
span of all its cycle effects. A cone is \emph{pointed} if it contains
no nonzero vector together with its negative. A component is
\emph{diagonal} when all its cycle effects are multiples of
\((1,1,1)\). A \emph{high} configuration has every coordinate above a
fixed enabling threshold. Thresholds and lattice conditions in the
following constructions depend only on the finite machine.

The following are proved supporting lemmas, rather than extra assumptions. Their detailed constructions and audits are retained in the accompanying research log and the named proof notes.

\begin{enumerate}[label=\textup{(\alph*)},leftmargin=*]

\item\label{item:ingredient-a} \textbf{The thin subsystem.} A fixed bounded-minimum region of a symmetric three-counter machine has an exact symmetric two-counter presentation. Store a chosen bounded physical coordinate and its value in finite control and keep the other two values as counters. Away from a region with two bounded coordinates, the chosen physical coordinate cannot change: a change within one update, whose size is bounded, forces two bounded coordinates. Tests of the stored coordinate are finite-control guards; any test of a free coordinate is a genuine scalar-axis visit. Finite change-of-marker gadgets at those visits preserve the exact physical endpoint values. Consequently \textbf{S2} supplies the reachability formulas needed for this fixed thin subsystem. Helper states representing totals or erasing counters are never mistaken for physical axes.

\item\label{item:ingredient-b} \textbf{Classifying update components.} In a strongly connected symmetric update graph, a non-diagonal pointed cycle cone is either positive or negative: all its nonzero cycle vectors have strictly positive total, or all have strictly negative total. A nonzero zero-total orbit or opposite total signs would create lineality. Non-diagonal nonpointed components have an exact
\emph{transfer-anchor} presentation. At selected anchor controls,
high vectors with the same total and appropriate coordinate residues
are mutually reachable. This is an equivalence at those anchors, not
at every control of the component. Actual ingress and egress relations
connect the anchors to other states. Its scalar port records the total
and residues. Diagonal low-to-low paths admit an actual uniformly thin replacement. Diagonal connectors incident to true high transfer ports have exact lower-dimensional presentations. They remain explicit hub/base connectors, not arbitrary relations inserted between productive cores.

\item\label{item:ingredient-c} \textbf{Productive face directions.} A positive pointed component is productive when its cone contains a vector with a negative coordinate. For every \(i\), the face diagonal \(d_i\) (zero in coordinate \(i\) and one in the other two) lies in the full interior of that cone. Indeed, from a negative-coordinate vector and its free swap obtain \((-a,b,b)\), where \(2b>a>0\). Writing
\(w_1=(-a,b,b)\), \(w_2=(b,-a,b)\), and \(w_3=(b,b,-a)\),
one can see the interior claim directly:
\[
 \frac{2b}{a}w_1+w_2+w_3
   =\frac{(2b-a)(a+b)}{a}(0,1,1).
\]
All three coefficients and the scalar on the right are positive, and
the three vectors are independent. Permutation gives the other two
face diagonals. The negative statement is reversed.

\item\label{item:ingredient-d} \textbf{Realizing high endpoints.} The high-endpoint cone/lattice lemma supplies a uniform threshold \(H\), a full-rank lattice \(\mathcal L_{\mathrm{lat}}\), and a fixed correction vector \(k\) such that suitable high endpoints \(u,v\) are joined whenever \(v-u\) is in the correct path lattice coset and \(v-u-k\) lies in the cycle cone. The proof uses all simple-cycle orbits, an affine-semigroup conductor, finite connector words, and bounded-error ordering of loops at a common control state. It does not assume that cycles at different states can be freely reordered. Boundary access is supplied by finite actual ingress/egress word catalogs: the access domains are upward Presburger sets by \textbf{S3}, hence have computable finite bases.

For a productive positive core, (c),(d) give a full-rank target period monoid on every bounded-coordinate face, with the balanced free direction in its interior. They also give input lifting: increasing the free source coordinates can be compensated by moving sufficiently far in that target face cone. Thus the marked productive face relation is weakly monotone on a sufficiently divisible grid. Its preimage of a grid-upward set is upward and effectively Presburger by \textbf{S3}. Repeated backward preimages stabilize by Dickson's lemma.

\item\label{item:ingredient-e} \textbf{Orthant components.} Call a positive component an \emph{orthant component} when all its cycles are coordinatewise nonnegative. In such a component, coordinate potentials show that any path ending with coordinate \(i\) bounded has that coordinate bounded throughout, by the endpoint bound plus a machine constant. The negative counterpart applies at the source. Across a finite condensation DAG the constants add. These orthant components provide thin replacements and the endpoint-box exceptions discussed below.

\item\label{item:ingredient-f} \textbf{Pumping finite scalar patterns.} Finite patterns of symmetric three-VASS pieces, exact symmetric two-CM pieces, and fixed tests/words have effectively semilinear relations by \textbf{S2} and \textbf{S3}. For prescribed scalar-axis endpoints they also have effective target pumping when target total minus source total is sufficiently large, and source pumping when it is sufficiently negative. The proof takes semilinear tuples of VASS schemes and scalar two-CM middle relations, chooses a positive-gain global period, and uses the water-filling offset

\begin{equation*}
w(x)=x-\proj_{\{z\geq 0:\,\sum_i z_i\leq\alpha\}}(x),
\end{equation*}


Here \(\proj\) is Euclidean projection onto the nonnegative
simplex of total at most \(\alpha\). Its argument \(x\) is a
nonnegative increment vector of the chosen semilinear period at a
point of the scheme, not the physical counter vector. If \(a\) is the
source increment, \(\alpha=\tot a\), and \(\beta e_i\) is the target
increment with \(\beta>\alpha\), then
\[
 w(a)=0,\qquad w(\beta e_i)=(\beta-\alpha)e_i.
\]
Thus the perturbation leaves the source fixed and adds a positive
axis increment at the target. VASS cycles are inserted in permitted orientations; full two-CM stages are handled through their scalar pumping theorem, without a false finite-flat-cover claim for counter machines. Bounded physical coordinates have homogeneous increment zero. The construction is effective by the semilinear presentations and finite witness enumeration.

For the positive patterns used below, gain filters and recorded high-core visits are preserved: inserted cycles have nonnegative total, so a fat gain cannot decrease below its cutoff; old high-core configurations survive with nonnegative offsets. Wide two-counter stages are translated by a fixed lower-bound vector before applying the argument. Reverse for negative patterns. This extension is not asserted for arbitrary filters on a full two-CM relation.

\end{enumerate}

\section{Why a bounded pair bounds large excursions}\label{sec:pair}

\subsection{The pair potential and the pumps it detects}

For any real vector, including a signed update, define
\[
 s_2(x)=\min\{x_1+x_2,x_1+x_3,x_2+x_3\}.
\]
For nonnegative \(x\), this is the sum of its two smallest coordinates.
The useful inequality is
\begin{equation}
 s_2(x+y)\ge s_2(x)+s_2(y).\label{eq:superadditive}
\end{equation}
Indeed every fixed pair sum of \(x+y\) is the corresponding pair sum
of \(x\) plus that of \(y\), and each summand is at least its own
minimum pair sum. Taking the minimum proves the inequality.

If \(\gamma=\delta_1\cdots\delta_k\) is a word, its \(a\)-weight is
\(\sum_{j=1}^k s_2(\delta_j)\). This need not equal
\(s_2(\sum_j\delta_j)\). The proof uses the sum of the individual
weights because it is additive along paths and can be analyzed by
finite graph potentials.

For an update \(\delta\) put \(t(\delta)=\sum_i\delta_i\) and let \(a(\delta)\) be the sum of its two smallest entries. If a positive symmetric component has a coordinate axis in the interior of its cone, there are actual independent axis-pump words at every state; expose them as full three-dimensional grid events.

Otherwise every cycle has nonnegative \(a\)-weight. Choose a shortest-path potential for \(a\). In the zero-reduced-\(a\) graph mark every vertex of an SCC supporting a positive-total cycle. To see why these are actual pumps, orient each update's two
smallest entries onto coordinates 1 and 2. A closed word of zero
\(a\)-weight and total \(T>0\) then has effect
\((r,-r,T)\) for some integer \(r\). Its global swap of coordinates
1 and 2 is also a closed word at the same control and has effect
\((-r,r,T)\). Concatenating them gives
\[
 (r,-r,T)+(-r,r,T)=(0,0,2T).
\]
A fixed buffer enables both words. Thus this is an actual axis-pump
word, not just an axis vector in a real cone. Full symmetry gives
pump words on the other two axes. Expose high visits as grid events; low visits belong to the fixed thin subsystem.

After deleting marked controls, every positive-total cycle has
strictly positive \(a\)-weight. There are only finitely many simple
cycles. Choose a positive rational \(\varepsilon\) no larger than
\begin{equation}
 \min_{\substack{\gamma\text{ a simple cycle}\\t(\gamma)>0}}
      \frac{a(\gamma)}{t(\gamma)}.\label{eq:epsilon}
\end{equation}
If there are no positive-total cycles, choose any positive
\(\varepsilon\); a finite path bound handles the remaining graph.
Zero-total cycles in the pointed component have zero effect; in the
residual graph their \(a\)-weight is zero as well. Consequently the
edge weight \(a-\varepsilon t\) has no negative cycle. A shortest-path
potential on the finite control graph bounds its weight on every path
below by a uniform \(-R\). Summing
\eqref{eq:superadditive} along the path gives

\begin{equation}
s_2(y)-s_2(x)\geq \varepsilon\bigl(\tot y-\tot x\bigr)-R, \label{eq:pair}
\end{equation}


where \(s_2\) is the sum of the two smallest coordinates. The same estimate holds for each fixed physical coordinate pair. It follows by summing \(a(\delta)-\varepsilon t(\delta)\), using

\begin{equation*}
s_2(x+\delta)\geq s_2(x)+a(\delta).
\end{equation*}


\subsection{From single excursions to a whole face phase}

Inequality~\eqref{eq:pair} controls one residual excursion. We still
have to interleave excursions with ordinary boundary corridors, which
may lower the pair potential. The next finite-graph argument controls
that loss.

Let \(V_0\) be the fixed outer face modes, including native thin-simulator modes but excluding later private word interiors. Keep the original fat-interface collar distinct from any larger finite marker values inside boundary catalogs or the thin simulator: enlarging the latter does not redefine the former or restart its construction. At mode \(v\) store a bounded marker \(b_v\) and a free pair \(z\). Put

\begin{equation*}
f_v(z)=b_v+\min(z_1,z_2).
\end{equation*}


Let \(B\) bound these outer marker values. Then
\[
 0\le f_v(z)-s_2(b_v,z_1,z_2)\le B.
\]
For example, if the marker is the smallest physical value, the two
expressions are equal. If both free values are smaller than the
marker, their difference is \(b_v-\max(z_1,z_2)\), which lies between
zero and \(B\). Thus \(f_v\) is a convenient finite-control version
of the physical pair potential. For an ordinary fixed-marker word from \(b\) to \(b'\) with free effect \((u,v)\), its guaranteed \(f\)-gain is

\begin{equation}
w_2=\min(u,v)+b'-b. \label{eq:ordinary-weight}
\end{equation}


Fix a large gain cutoff \(\Co\). Only \emph{positive} same-marker pieces of gain at most \(\Co\) are compiled into the positive ordinary graph; use a separate reversed construction for negative pieces. Nonnegative reduced total costs bound their word lengths, after deleting zero-effect cycles. If a high-core visit must be retained, split there first and shorten each half separately. All lower total and pair losses are bounded independently of \(\Co\). A bounded-gain cross-marker piece has uniformly bounded coordinate displacement by the same reduced-cost estimate: the source marker and the target marker are both bounded at both endpoints. Choose \(W\) later above this bound and treat these pieces as scalar-axis endpoint patterns. Thus every remaining residual fat edge really has gain greater than \(\Co\).

Every finite word catalog is closed under \emph{all} independent free-coordinate swaps at its individual steps. Each complete oriented word is an outer edge. This is necessary for an exact symmetric two-VASS chain presentation; whole-word swaps alone would be insufficient. All expansions retain the physical total and marker trace, and \eqref{eq:pair} supplies the same cutoff-independent pair-loss bound.

An ordinary closed word can be oriented to have free effect \((A,T-A)\), where \(A\) is its accumulated \(w_2\) and \(T\) its total. A negative pair-weight cycle cannot simply be ignored. Its
actual words identify the exceptional behavior that must first be
exposed or normalized:
\begin{itemize}[leftmargin=*]
\item If \(A<0\) and \(T<0\), concatenate the free-coordinate swap
to obtain effect \((T,T)\), a negative free diagonal.
\item If \(A<0\) and \(T=0\), the effects \((A,-A)\) and
\((-A,A)\) give an actual transfer mechanism.
\item If \(A<0<T\), write the effect as \((-a,b)\), so \(b>a>0\).
If \(W\) is that closed word and \(\overline W\) its swap, then
\(W^a\overline W^{\,b}\) has effect \((b^2-a^2,0)\).
Swapping gives the other positive axis pump.
\item If \(A=0<T\), the cycle already has effect \((0,T)\),
and its swap pumps the other coordinate.
\end{itemize} All have finite actual word catalogs and enabling buffers.

After these marked cases are dealt with by Section~\ref{sec:normal}, surviving ordinary cycles have nonnegative pair weight. Let \(-P_2\) bound each ordinary edge's weight. Choose \(\Co\) so that

\begin{equation}
\begin{aligned}\Gamma&=\varepsilon\Co-R-B,\\ \Gamma-1&>(|V_0|-1)P_2.\end{aligned}\label{eq:cutoff}
\end{equation}


Give each fat edge formal weight \(\Gamma-1\), and each ordinary
edge its guaranteed pair weight. A cycle with no fat edge is
nonnegative. In a cycle with fat edges, erase ordinary subcycles
between consecutive fat edges. Each remaining ordinary subpath has
at most \(|V_0|-1\) edges and loses at most \((|V_0|-1)P_2\).
Condition~\eqref{eq:cutoff} pays for that loss with each fat edge.
Hence this formal graph has no negative cycle.

A control potential \(h_2\) therefore makes
\[
 F(v,z)=f_v(z)+h_2(v)
\]
nondecreasing on ordinary edges and increasing by at least one on
every fat edge: the latter's actual gain is at least \(\Gamma\), one
more than the formal weight used to construct the potential.

\begin{lemma}[The excursion-count bound]\label{lem:count}
In a residual positive phase whose target satisfies \(f_v\le M\),
the number of complete fat edges is at most
\[
 M+\max_{v\in V_0}h_2(v)-\min_{v\in V_0}h_2(v).
\]
This bound is uniform in the unrestricted coordinates of both endpoints.
A negative phase has the corresponding bound from its source pair.
\end{lemma}
\begin{proof}
If there are \(k\) fat edges, monotonicity of \(F\) gives
\(k\le F(\text{target})-F(\text{source})\).
Use \(f\ge0\) at the source and \(f\le M\) at the target.
Reverse the machine and the endpoint roles for the negative statement.
\end{proof}
This is the step that turns a geometric statement about a pair into a
finite enumeration of run patterns. It does not bound total height:
the third coordinate can still grow without bound.

\subsection{Why the constants are not circular}

Only after \(\Co\) and the finite catalogs are fixed choose the wide/axis threshold \(W\). If an axis cut lies inside a private word, keep its word-position and mandatory-continuation tag. Such scalar ports are endpoint adapters, not new vertices of \(V_0\). A fixed maximum prefix displacement \(A\) changes the endpoint pair bound to at most \(B+W+A\). Thus these cuts do not create circular constants. A later scalar pumping cutoff \(\Cs\) and common modulus \(D\) do not change \(\Co\) or \(W\). Potentials on \(V_0\) lift unchanged to residue refinements. Hub canonical bounded coordinates retain the \emph{original} transfer modulus; refining \(D\) changes the scalar-total residue classes, not the sizes of those bounded representatives. Repeating actual nonnegative or nonpositive coordinate pump words to obtain \(D\)-multiple effects needs no larger enabling threshold when both event endpoints remain eligible.

\section{Normalizing sign changes and transfer visits}\label{sec:normal}

The pair estimate applies inside a phase of one sign. A general run
can change sign, visit transfer components, and cross axes. This
section explains how to cut such a run into the pieces that the oracle
will handle.

The eventual output must satisfy three separate requirements:
\begin{enumerate}[leftmargin=*]
\item every normalized path between original endpoints represents an
actual machine run (soundness);
\item every actual run has a normalized representative (completeness);
\item all connectors use \emph{one fixed relation} \(K\), with an
algorithm deciding its box queries (effectivity).
\end{enumerate}
A termination measure for rewriting proves only part of the second
requirement. The exact macro realizations, interface rules, and oracle
below are also needed.

\subsection{Peaks, valleys, and a decreasing measure}

We now describe the exact normalization rather than assuming that finitely many semilinear macro relations have a decidable closure.

Productive high occurrences have signs \(P,N\). Between physical scalar-axis cuts, every actual thin corridor has a fixed bounded physical marker. A \(P\)-to-\(N\) \emph{peak} first gains and then loses total.
The intermediate height can be made large enough to connect compatible
endpoints. Its relation across a fixed-marker corridor is an exact
finite union of Cartesian products
\[
 \bigcup_{\ell=1}^r E_\ell\times F_\ell.
\]
The assertion is that \emph{every} pair in each displayed product has
a realizing run, not merely that both projections are reachable:
 finite residue ingress/egress catalogs and a positive face-interior cone followed by the opposite diagonal pump connect every pair in the chosen product. Expose each product as a zero-dimensional event.

An \(N\)-to-\(P\) \emph{valley} first loses and then gains total.
Here the idea is to lower the middle of the run while preserving the
tests, until it approaches an axis. The cores on either side absorb
the changed boundary values. Thus the valley has an actual replacement
with a physical axis cut. Select the \emph{last} negative productive SCC with an actual high visit and the \emph{first} positive productive SCC with such a visit. Freeze everything outside these two core occurrences. Shallow intervening productive pieces are thin; intervening nonproductive pieces have orthant coordinate potentials. Outside an existing axis cut their bounded-marker identity is fixed, say \(i\). If \(M\) is the least free-coordinate value over the genuine thin middle, all adjacent orthant tails have free values at least \(M-R_0\). Subtract a sufficiently divisible multiple of \(d_i\) from this middle until one free coordinate is bounded. Its marker and all its tests are unchanged. Reconstruct only the selected single-SCC core pieces: their facet deficits are uniformly bounded, and the signed interior directions \(-d_i,+d_i\) repair them. Low-marker SCC exits use finite boundary catalogs. No claim is made that a whole multi-SCC displacement belongs to one core's cone.

The same valley/peak replacements hold with one side replaced by an available ordinary free-diagonal loop. These are the replacements needed for ordinary contrary-sign cycles.

Termination is explicit. Before choosing \(\Co\) label the original finite list of deep productive occurrences. A peak consumes origins; a valley separates unlike adjacent origins by an axis cut, retaining their signs. The lexicographic pair (retained origins, unlike-sign adjacencies in uncut intervals) decreases. Reconstructed signed pieces are atomic relations in their \emph{same} residual signed graphs; they are not recursively expanded at every new physical collar visit.

After the sign-specific small-word catalogs are fixed, label ordinary adverse wide visits. In a positive interval, a negative-diagonal visit after a fat is removed with the \emph{last} preceding fat by a one-fat peak, decreasing the number of fat origins. If all such visits precede the first fat, choose the \emph{last} one and perform a one-fat valley; the selected visit and inserted negative loops lie in the base-only prefix before the new axis cut. Thus (retained fat origins, original adverse visits in still-fat-containing intervals) strictly decreases. The negative construction is reversed. New private words are confined to base-only pieces, atomic signed relations, or mandatory adapters; they are not added as unrestricted edges of the residual ordinary graph.

\subsection{Transfer hubs and mandatory private events}\label{sec:hubs}

At a genuine transfer anchor, compatible high vectors of a fixed total
are mutually reachable. A scalar \emph{hub} records this total and
residue information and uses a fixed canonical representative with two
bounded coordinates. Exact adapters connect it to the physical high
representatives.

This compression can hide a large pair. For example, a representative
near an axis can encode the same large total as a balanced vector with
three large coordinates. Therefore an arbitrary expansion from a hub
cannot be treated as beginning with a bounded physical pair.
The next two mandatory event rules prevent that mistake.

Grid visits are exposed as actual independent-pump events. A genuine high transfer visit \emph{after} a positive fat is handled by its last-fat-to-hub relation \(R\) followed by a private scalar \(\Up\). Here \(R\) is PA by \textbf{S2} and \textbf{S3} and \(R;\Up=R\): pump a positive face period at the last core and translate the following fixed-marker corridor. The reverse is a private \(\Down\) \emph{before} a hub-to-negative-fat relation. In symbols, the two relevant factorizations have the form
\[
 \text{last positive core}\ ;R;\Up;\text{hub},
 \qquad
 \text{hub};\Down;R;\text{first negative core}.
\]
Here \(R;\Up=R\) in the first case and \(\Down;R=R\) in the
second. A private \(\Up\) is sound because every entrance passes
through its specified \(R\); it is not asserted to be a physical
unconditional loop. The dual restriction holds for \(\Down\).
In particular, a negative scalar phase cannot begin by freely expanding
a canonical hub into arbitrary high representatives and then use a
fictitious bounded source pair.

A hub before the first positive fat may be represented by a base-only detour to its canonical scalar representative and the exact inverse hub adapter. The latter exposes every high representative with the same total/residues. The positive residual pair count begins after this adapter; no earlier high pair is bounded by its canonical representative. Reverse for a negative phase. The same original-visit measure terminates these hub cuts.

Initial/final productive occurrences are included in the sign normalization, with exterior endpoints frozen. Orthant endpoint exceptions are retained explicitly: a negative-orthant initial fragment maps a source box to an upward face set or a PA bounded-coordinate strip; a positive-orthant initial fragment ending on a face is uniformly thin. Reverse for terminal fragments. An arbitrary mixed-sign \textbf{S3} relation is \emph{not} silently prepended to a negative phase.

\subsection{The fixed connector relation and its contract}

A connector is built from exact lower-dimensional corridors, normalized
residual phases, and the prescribed endpoint adapters. It can pass
through arbitrarily many physical scalar-axis interfaces. It does not
pass freely through a designated independent-pump or reset event:
those are the events separated out for the outer saturation.

All modes and relation types are fixed from the machine: physical faces, full three-dimensional pump interfaces, physical scalar axes (including private-word continuation tags), true total/residue hubs, Cartesian reset ports, and mandatory private scalar pump ports. Erase origin labels after the termination proof. Let \(K\) be the resulting fixed connector language between designated events, allowing arbitrary composition only through its scalar-axis interfaces and exact lower subsystem. This is a fixed sound normalized relation, not a query-dependent family and not the unrestricted reachability relation in disguise.

Soundness follows from actual word/macro realizations, true hub equivalence, Cartesian reset soundness, and absorption of private \(R;\Up\) or \(\Down;R\). Completeness is the terminating normalization just proved. Consequently the normalized reachability language is
\begin{equation}
 L=K;\bigl((U\cup V);K\bigr)^*.\label{eq:language}
\end{equation}
Here \(U\) collects the exposed upward events and \(V\) the downward
events, including the private events with their mandatory adapters.
Original reachability between original endpoint charts equals this
relation. Importantly, \(K\) need not be transitive and need not have a
semilinear full relation. The next section supplies exactly the oracle
that the final saturation uses.

\section[The box oracle for K]{The box oracle for \(K\)}\label{sec:oracle}

The normal form is useful only if its connector can be queried without
already solving the original reachability problem. We now construct
that oracle. It answers a Boolean question
\[
 K(S,T)\quad\Longleftrightarrow\quad
       \exists x\in S\ \exists y\in T:\ xKy.
\]
It is not required to return a formula for all pairs in \(K\).

Fix a positive modulus \(D\), ultimately chosen divisible by the
finitely many structural pumping periods. In a chart \(\N^d\), a
\emph{grid box} is a product of intervals, each in a fixed residue
modulo \(D\). Each interval has a finite lower bound and either a finite
or infinite upper bound. Finite unions of such boxes form an effective
Boolean algebra \(\mathcal B\). The order within a fixed chart and
residue is componentwise increase by a vector in \(D\N^d\).

\subsection{The two-dimensional box dichotomy}

The reason for using boxes is a simple exhaustive alternative. A box
in \(\N^2\) either has both coordinates unbounded above, or has at
least one bounded coordinate. For example, ignoring residues,
\[
 [a,\infty)\times[b,\infty)
 \qquad\text{and}\qquad
 [a,A]\times[b,\infty)
\]
illustrate the two cases. On a face with bounded marker, the second
case bounds a physical pair. These two cases are handled by different
arguments.

First consider one positive residual face-to-face phase.

\paragraph{Both free target coordinates are unbounded.}
The target is grid-upward. The marked productive face relations have
the following weak monotonicity property on a sufficiently divisible
grid:
\begin{equation}
 xRy,\quad x'\ge_D x
 \quad\Longrightarrow\quad
 \exists y'\ge_D y:\ x'Ry'.\label{eq:weak-monotonicity}
\end{equation}
Here source and target modes are unchanged by the comparisons.
The geometric reason is that an arbitrary free source increment can
be compensated by a sufficiently large balanced increment on the
target face. The latter direction is interior to the productive cone;
the high-endpoint lemma realizes the resulting splice. Recorded
high-core visits and the large-gain filter are retained by this
construction. Ordinary fixed-marker corridors have the corresponding
translation property.

Therefore the preimage of an upward set is upward. Since each macro
relation is PA, one can repeatedly compute the PA predecessors of the
target and take their union. The increasing sequence of upward sets
stabilizes by Dickson's lemma, and PA equality detects stabilization.
An arbitrary PA source set can be intersected with the resulting PA
predecessor. The target's upwardness, not the source's, is what makes
this argument work.

\paragraph{At least one free target coordinate is bounded.}
Together with the stored marker, that bound gives \(f_v\le M\) for
a computable \(M\). Lemma~\ref{lem:count} bounds the number of fat
edges by a computable \(L(M)\). Thus only finitely many patterns of
the form
\begin{equation}
 O^*;F_1;O^*;F_2;\cdots;F_k;O^*,
 \qquad 0\le k\le L(M),\label{eq:finite-patterns}
\end{equation}
need be considered. Here \(O^*\) is an ordinary corridor relation,
and each \(F_i\) is one of the finitely many fat macro types. The
corridors are exact symmetric two-VASS relations, the fat macros are
\textbf{S3} relations, and the adapters are fixed. \textbf{S2},
\textbf{S3}, and PA composition compute each finite pattern and test
it against the endpoint constraints.

The pattern bound may depend on the target box. This does not change
\(K\): it changes only how much of its fixed grammar must be examined
for that particular query. Negative phases are handled by reversing
the whole phase and its endpoint constraints.

\begin{remark}
The dichotomy fails for arbitrary PA targets. For example,
\(\{(n,n):n\in\N\}\) is neither grid-upward nor contained in a
bounded-coordinate strip. Accordingly, the argument above allows
arbitrary PA sources but uses boxes, upward sets, or bounded-pair
strips as targets. This is why the completion theorem asks for a box
oracle rather than an oracle for arbitrary PA endpoint sets.
\end{remark}

\subsection{Endpoints at full charts, hubs, and private adapters}

For full three-dimensional target boxes, a fully unbounded box gives an upward ingress domain by translation. If exactly one coordinate \(i\) is bounded, split its finite interval. At a productive residual core, \(d_i\) is interior to its cone: translate the prefix by an arbitrary source increment and the suffix by a sufficiently large multiple of \(d_i\), then splice at the high core using the cone/lattice lemma. This yields an upward PA face source domain for that branch. In the orthant branch a coordinate potential bounds the corresponding source coordinate; if it is free, this is a bounded-pair PA strip, and if it is the marker, free translation gives upwardness. Two bounded target coordinates give a bounded-pair preimage by the residual pair inequality. \textbf{S3} computes the exact ingress relations; no semantically punctured graph is substituted for a symmetric VASS. Reverse all signs and endpoint roles for negative phases.

For a last-positive-fat hub adapter \(R\), an unbounded scalar target tail has an upward PA source domain: use the cofinal positive face periods to enable a finite residue-compatible corridor word, then pump beyond the target threshold. A bounded scalar target has finitely many possible source vectors: a finite control potential for total weight gives \(\tot{\text{source}}\le \tot{\text{target}}+c\) for a fixed constant \(c\), so bounded target total bounds all source coordinates. Thus the preceding positive phase is queried against an upward set or a bounded-pair set. A scalar hub source tail denotes an upward PA family of high representatives; a bounded scalar source denotes finitely many representatives. The negative cases are dual and obey the mandatory-event grammar in Section~\ref{sec:normal}.

Private reset adapters are finite unions of grid-upward products. Fixed partial-word adapters preserve boxes by translations and finite enabling bounds. PA bounded-coordinate strips are handled by bounded-pattern enumeration; alternatively they decompose into finitely many unary semilinear slices. Internal calculations may use a query-dependent modulus \(M\) divisible by the fixed structural periods and \(D\). They need not force every resulting PA congruence to divide the \emph{outer} saturation period \(D\): \(K\) returns only a Boolean, and its internal refinement changes no transitions or thresholds. The outer upward sets are \(D\)-upward because of the actual events, independently of these internal calculations.

\subsection{First and last axis visits: the exact reduction}

So far we can handle normalized passages without an intervening chosen
axis port. There can still be arbitrarily many \emph{internal} physical
axis visits. The following decomposition isolates precisely that
remaining closure problem. Compute the unary PA first-axis set \(I\) reachable from the source boxes without an intervening chosen scalar-axis port, and the unary PA last-axis set \(J\) from which the target boxes are reachable without another one. The preceding cases compute these sets: bounded pairs give finite PA patterns, while upward cases give effective PA predecessors.

Write \(I=I(S)\) and \(J=J(T)\) to display their dependence on the
endpoint boxes. Empty prefixes or suffixes are allowed when an endpoint
already lies at a scalar interface. These are formula computations,
not merely Boolean queries: in a bounded-pair branch, keep the
remaining scalar free in the finite-pattern PA formula; in an upward
branch, compute the PA predecessor and restrict it to the scalar
interface. Use the reversed construction for the corresponding
first-axis cases.

Here an axis visit is a selected scalar interface in the normalized
grammar. An atomic signed macro may have a realizing run with
additional physical collar visits; it is not recursively cut at all
of them. This convention is the same one used in the normalization
measure.

Let \(Z(S,T)\) denote the already decidable no-axis case. Let
\(E\) denote the residual scalar passages to be described below, and
\(L_{\mathrm{low}}\) the exact lower-subsystem relation on the same
scalar interfaces. The full connector query is exactly
\begin{equation}
 \boxed{\quad
 K(S,T)\ \Longleftrightarrow\ Z(S,T)\ \lor\
 \bigl(I(S);(E\cup L_{\mathrm{low}})^*\bigr)\cap J(T)\ne\varnothing.
 \quad}\label{eq:kernel-reduction}
\end{equation}
Every path either has no axis visit, or has a first visit in \(I\),
some sequence of scalar passages, and a last visit in \(J\).
Conversely each such factorization concatenates exact connector
pieces. This explains how the local pair bound is connected to the
Boolean oracle for the whole fixed relation \(K\).

\subsection{Scalar passages and bounded-band compilation}

\paragraph{The remaining problem is one-dimensional.}
There can be arbitrarily many visits to scalar-axis ports, even though
each individual passage between such ports has a finite description.
We therefore separate the computation of one passage from the closure
under arbitrarily many passages.  Write
\[
X_{\mathrm{ax}}=Q_{\mathrm{ax}}\times\N
\]
for the finite collection of scalar ports, with their port modes included
in the finite set \(Q_{\mathrm{ax}}\).  At a physical axis port the scalar
records the one unbounded coordinate, together with a fixed offset if
that is part of the chosen port convention.  The bounded coordinates and
all continuation tags are recorded in the mode.  Canonical hub ports
also have their already specified scalar encoding.

Let \(E\) be the relation on this space supplied by the residual
positive and negative passages between consecutive scalar-axis ports.
The mandatory hub rules from Section~\ref{sec:hubs} remain part of this
definition. A positive passage admitted here has a genuine physical
axis port at its terminal end. An initial hub expansion is an endpoint
adapter, and the pair estimate starts after that adapter. Dually, a
negative passage has a genuine bounded physical pair at its source.
A last-positive-fat-to-hub passage, and its reversed negative case,
cross a mandatory private event and are handled by the endpoint cases
above; they are not inserted freely into \(E\).
Base-only hub passages belong to \(L_{\mathrm{low}}\).

The important point is that \(E\), before taking its transitive closure,
is effectively Presburger-definable.  Indeed, a positive passage has a
bounded physical pair at its terminal port, so the pair estimate gives
a uniform bound on its number of fat pieces.  A negative passage has
the dual bound from its initial port.  There are therefore only finitely
many patterns to consider.  The ordinary corridors in each pattern have
exact lower-dimensional presentations, its fat pieces have the
\textbf{S3} presentations, and its adapters are fixed relations.
\textbf{S2}, \textbf{S3}, and Presburger composition compute their union.
The bound here depends on the fixed port conventions, not on the
unbounded scalar values.

There is a second property of \(E\), independent of its Presburger
description.  Exact fixed-value zeroizers turn the collar endpoints of
these finitely many patterns into literal axis endpoints, so that the
finite-pattern axis-pumping result of
Section~\ref{sec:ingredients}(f) applies.  Taking a common period and
enlarging the cutoff to absorb the fixed port offsets gives constants
\(D>0\) and \(\Cs\geq0\) with the following implications, for all port
modes \(p,q\):
\begin{align}
 (p,n)E(q,m),\quad m-n>\Cs
 &\quad\Longrightarrow\quad (p,n)E(q,m+D),
 \label{eq:scalar-positive-pump}\\
 (p,n)E(q,m),\quad n-m>\Cs
 &\quad\Longrightarrow\quad (p,n+D)E(q,m).
 \label{eq:scalar-negative-pump}
\end{align}
These are geometric pumping statements; Presburger-definability alone
would not imply them.  The zeroizers are auxiliary proof gadgets, and
their internal heights do not change the fixed collar width \(W\).

\paragraph{Small scalar changes belong to the two-counter part.}
Partition \(E\) according to its scalar difference:
\[
\begin{aligned}
 E_+&=\{((p,n),(q,m))\in E:m-n>\Cs\},\\
 E_-&=\{((p,n),(q,m))\in E:n-m>\Cs\},\\
 E_0&=\{((p,n),(q,m))\in E:|m-n|\leq\Cs\}.
\end{aligned}
\]
All three relations are effective PA relations.  In the bounded band
\(E_0\), only finitely many differences \(d=m-n\) occur.  For every
choice of \(p,q,d\), the legal starting scalars form the unary PA set
\[
 G_{p,q,d}=\{n\in\N:((p,n),(q,n+d))\in E_0\}.
\]
Thus an \(E_0\)-step is a fixed scalar shift, enabled by a unary
Presburger guard.  This restricted form is what makes it possible to
add these steps to the exact lower symmetric two-CM subsystem.

Here are the encoding details.  At port \(p\), represent the scalar
\(n\) by a two-counter configuration \((x,0)\) with
\(n=x+c_p\), where \(c_p\) is the fixed offset of that port.  If a port
uses the opposite active-counter orientation, the usual exact transfer
adapter puts it in this orientation before the guard is checked.
A unary PA set is a finite union of exceptional values and
threshold-and-residue tails.  Choose a common modulus \(M\) for all
the guards and track the encoded total \(x+y\bmod M\) in finite
control throughout the lower subsystem.  Permuting a two-counter
update does not change its total, so this refinement preserves
symmetry.  At a port, where \(y=0\), it determines \(x\bmod M\)
and hence \(n\bmod M\).  The finitely many threshold and exceptional
value tests are implemented by fixed decrement/restore words and zero
tests.  Finally, a shift \(m=n+d\) becomes
\[
 x'=x+d+c_p-c_q.
\]
The unused counter is tested zero in the shift gadget to enforce the
required port orientation.  These are finite gadgets attached to the
exact lower subsystem; the construction does not compile an arbitrary
binary PA relation into two counters.

Let \(L_{\mathrm{low}}\) denote the relation furnished by that lower
subsystem on its scalar interfaces, and put
\begin{equation}
 \Lambda=(L_{\mathrm{low}}\cup E_0)^*.
 \label{eq:scalar-lower-closure}
\end{equation}
The combined two-CM presentation and \textbf{S2} compute \(\Lambda\)
as a PA relation.  In particular, \(\Lambda\) is reflexive and
transitive.  What remains is to decide
\begin{equation}
 \bigl(I;(\Lambda\cup E_+\cup E_-)^*\bigr)\cap J
 \neq\varnothing
 \label{eq:scalar-query}
\end{equation}
for the unary PA first-axis and last-axis sets \(I,J\) obtained above.

\subsection{The scalar backward and forward saturations}
Call a subset \(Y\subseteq X_{\mathrm{ax}}\) \(D\)-upward if
\((q,n)\in Y\) implies \((q,n+D)\in Y\).  In each mode and residue
class, such a set is either empty or a tail.  Hence it has a finite
basis, and an increasing sequence of such sets must stabilize.  By
\eqref{eq:scalar-positive-pump}, every image under \(E_+\) is
\(D\)-upward.  By \eqref{eq:scalar-negative-pump}, every preimage
under \(E_-\) is \(D\)-upward.  These statements hold for arbitrary
sets of starting or ending scalars.

First allow any number of negative steps, each followed by a lower
passage.  Write
\[
 \mathcal R_-=(E_-;\Lambda)^*.
\]
For a unary PA target set \(T\), compute
\begin{equation}
 T_0=T,\qquad
 T_{k+1}=T\cup E_-^{-1}\bigl(\Lambda^{-1}(T_k)\bigr).
 \label{eq:scalar-backward-saturation}
\end{equation}
Every operation in this formula is a PA operation on known PA
relations and sets.  The sequence is increasing, and after the fixed
part \(T\) is removed its growing part is \(D\)-upward.  It therefore
stabilizes.  PA equality detects the first fixed point, which is exactly
\(\mathcal R_-^{-1}(T)\).  We have computed this preimage for the
given target; we have not asserted a PA description of the full binary
relation \(\mathcal R_-\).

Next, for a unary PA source set \(S\), consider the set
\[
 F(S)=S;\Lambda;\mathcal R_-;E_+.
\]
It is \(D\)-upward because its last step is positive.  Although we do
not have the full relation \(\mathcal R_-\), we can decide whether
\(F(S)\) meets any unary PA test set \(H\):
\begin{equation}
 F(S)\cap H\neq\varnothing
 \quad\Longleftrightarrow\quad
 (S;\Lambda)\cap
 \mathcal R_-^{-1}\bigl(E_+^{-1}(H)\bigr)\neq\varnothing.
 \label{eq:scalar-forward-oracle}
\end{equation}
The right side is effective: first compute the PA set
\(E_+^{-1}(H)\), then apply
\eqref{eq:scalar-backward-saturation}, and finally use PA emptiness.

This intersection test computes a finite basis of \(F(S)\).
Start with the empty upward approximation \(A\).  Test whether
\(F(S)\) meets its PA complement.  If so, enumerate singleton
configurations until a witness \((q,n)\in F(S)\setminus A\) is
found, and add the whole ray \(\{(q,n+kD):k\in\N\}\) to \(A\).
The approximation remains a subset of \(F(S)\), and each enlargement
is strict.  The ascending-chain condition gives termination, while
\eqref{eq:scalar-forward-oracle} gives the stopping test.  Consequently
\(F(S)\) is effectively obtained as a PA set.

Now iterate
\begin{equation}
 S_0=I,\qquad S_{k+1}=I\cup F(S_k).
 \label{eq:scalar-forward-saturation}
\end{equation}
Again the growing part is \(D\)-upward, so this is an effective
terminating iteration.  Its limit is \(I;\mathcal A\), where
\[
 \mathcal A=(\Lambda;\mathcal R_-;E_+)^*.
\]
To see that these two computations suffice, split any word over
\(\Lambda,E_+,E_-\) immediately after each positive step.  Before
the next positive step, and after the last one, only lower and negative
steps remain.  Reflexivity and transitivity of \(\Lambda\) therefore
give the exact factorization
\begin{equation}
 (\Lambda\cup E_+\cup E_-)^*
   =\mathcal A;\Lambda;\mathcal R_-.
 \label{eq:scalar-factorization}
\end{equation}
The original scalar question \eqref{eq:scalar-query} is consequently
equivalent to the PA emptiness test
\begin{equation}
 (I;\mathcal A;\Lambda)\cap\mathcal R_-^{-1}(J)
 \neq\varnothing.
 \label{eq:scalar-final-test}
\end{equation}
The first set is computed by the forward saturation followed by one
\(\Lambda\)-image, and the second by the backward saturation.
This proves the required scalar set-to-set decision procedure.  It
does not claim that the entire scalar closure is an effectively
semilinear binary relation, or that all its unary images and preimages
are computed.


Together with the no-axis case,
\eqref{eq:scalar-final-test} decides \eqref{eq:kernel-reduction}, and
hence \(K(S,T)\) for every pair of boxes. The oracle is now complete:
its only unbounded iterations are the upward saturations just justified.
Its query-dependent finite bounds never redefine \(K\).

\section{Adding the independent pumps back}\label{sec:completion}

We have obtained a fixed connector \(K\) and its box oracle. The final
step is independent of the geometry: it uses only that oracle and
independent coordinate pumps. This section states the general argument
so that its precise algorithmic requirements are visible.

Let \(X\) be the finite disjoint union of the interface charts
\(\{q\}\times\N^{d(q)}\), with \(d(q)\le3\). Use the grid boxes,
the algebra \(\mathcal B\), and the fixed period \(D\) from
Section~\ref{sec:oracle}.

Let \(K\) be a \emph{fixed} relation on \(X\). It need not be transitive or semilinear. Suppose the question

\begin{equation}
K(S,T)\quad\Longleftrightarrow\quad\exists x\in S\ \exists y\in T:\ xKy. \label{eq:box-query}
\end{equation}


is decidable for every \(S,T\) in \(\mathcal B\).

At selected eligible charts let \(U\) independently add arbitrary multiples of \(D\) to the coordinates. Eligibility is a grid upset given by an effective finite basis. At other selected charts let \(V\) independently subtract multiples of \(D\), keeping both endpoints eligible. Both events allow zero addition/subtraction. An event may change an input chart to a designated output chart of the same dimension. Dimension zero is allowed.
Independence means, for example, that on a two-counter chart an upward
event can send \((x,y)\) to \((x+kD,y+\ell D)\) with \(k,\ell\)
chosen separately. A diagonal event allowing only \(k=\ell\) would
not supply the product-order upward closure used below.

\begin{theorem}[Grid-pump completion]\label{thm:completion}

Then the relation \(L\) in \eqref{eq:language} has decidable box-to-box emptiness.

\end{theorem}

\begin{proof}
\textbf{Step 1: make Dickson's lemma effective.}

Box sets are effectively closed under Boolean operations, \(U\)-preimages and \(V\)-images. A grid-upward set has a finite basis by Dickson's lemma. Moreover an intersection oracle for an upward set \(Y\) against all boxes computes its basis: maintain a finite upward approximation \(A\); test whether \(Y\cap(X\setminus A)\) is empty; if it is nonempty, enumerate singleton witnesses outside \(A\) and add one generated upset. Each enlargement is strict, so the ascending-chain condition for upward subsets of a well-quasi-order terminates the procedure. The oracle gives the effective stopping test.

\medskip
\noindent\textbf{Step 2: close backwards under negative events.}

Put \(\mathcal C=(V;K)^*\). This permits any number of downward
events, each followed by a connector. For a box set \(T\), iterate

\begin{equation}
\begin{aligned}T_0&=T,\\ T_{n+1}&=T\cup V^{-1}\bigl(K^{-1}(T_n)\bigr).\end{aligned}\label{eq:backward}
\end{equation}


The added part
\(Y_n=V^{-1}(K^{-1}(T_n))\) is upward: if a configuration can first
subtract to reach a connector source, any larger eligible configuration
can subtract to that same source. Its intersection with a box \(E\) is decided by the available query

\begin{equation*}
K\bigl(V(E),T_n\bigr).
\end{equation*}


Thus \(Y_n\) is effectively computed by the basis procedure. The sequence consists of the fixed \(T\) together with an increasing upward part, so it stabilizes effectively at \(\mathcal C^{-1}(T)\).

\medskip
\noindent\textbf{Step 3: close forwards under positive events.}

For a box set \(S\), put \(f(S)=S;K;\mathcal C;U\). This means:
take a connector, allow negative-event closure, and then take one
upward event. The result is upward because of its final event. Its
intersection with a test box \(E\) is decided by

\begin{equation}
f(S)\cap E\neq\varnothing\quad\Longleftrightarrow\quad K\bigl(S,\mathcal C^{-1}(U^{-1}(E))\bigr). \label{eq:forward-query}
\end{equation}


Again the basis procedure computes \(f(S)\). Iterating \(S_{n+1}=S_0\cup f(S_n)\) stabilizes at \(S_0;\mathcal A\), where \(\mathcal A=(K;\mathcal C;U)^*\). Finally, split any word in \(L\) immediately after each upward
event. All intervening event blocks consist only of negative events.
This gives

\begin{equation}
L=\mathcal A;K;\mathcal C. \label{eq:factorization}
\end{equation}


The repeated blocks end in \(U\); the remaining suffix is
\(K;\mathcal C\). This identity does not use transitivity of \(K\).
Both \(S_0;\mathcal A\) and \(\mathcal C^{-1}(T)\) have been
computed as box sets, so the final question is the available query
\[
 K\bigl(S_0;\mathcal A,\mathcal C^{-1}(T)\bigr).
\]
This proves the theorem.

\end{proof}

\paragraph{Private interfaces.}

The private interfaces used in the construction fit the theorem as follows. A Cartesian relation \(E\times F\) can be exposed as an event at a zero-dimensional private port: check \(E\) before the event and generate \(F\) after it. A right-saturated relation \(R\) with \(R;U=R\) can be factored through a private \(U\) event, provided \emph{every} entrance into that event passes through \(R\). Soundness then follows by absorbing \(R;U\) into \(R\). Dually use \(V;R=R\). Such a private event is not asserted to be a physical self-loop.


\section{Assembling the decision algorithm}\label{sec:conclusion}

Given an original machine and two configurations, the algorithm is now
specified in the following order.
\begin{enumerate}[leftmargin=*]
\item Construct the finite interface types, the thin presentations,
the signed core relations, and their actual ingress/egress adapters.
Expose the initial pump and transfer cases, and prepare the
peak/valley normalization rules and their decreasing measures.
\item Establish the residual pair estimates on the fixed outer graph.
Choose \(\Co\), then the finite word catalogs and axis width \(W\).
Complete the prescribed normalization and instantiate the fixed
connector grammar. Use its scalar patterns to choose \(\Cs\) and a common period
\(D\). These later choices do not change the earlier outer graph or
its potential.
\item Implement \(K(S,T)\) by the two-dimensional box dichotomy,
endpoint adapters, and the first-axis/last-axis reduction
\eqref{eq:kernel-reduction}. Its scalar closure is decided by
\eqref{eq:scalar-backward-saturation}--\eqref{eq:scalar-final-test}.
\item Apply Theorem~\ref{thm:completion} to the singleton source and
target boxes. Its backward negative closure and forward positive
closure terminate and return the desired yes/no answer.
\end{enumerate}

The two assumed semilinearity algorithms have different concrete roles.
\textbf{S3} supplies formulas for actual update-only macro pieces and
their constrained endpoint relations. \textbf{S2} supplies formulas
for the thin subsystem and, after the bounded-band compiler, for the
combined lower scalar closure \(\Lambda\). Neither assumption is
being applied to unrestricted three-counter zero-test reachability.

Section~\ref{sec:normal} supplies a finite exact event/kernel presentation of the original machine. Section~\ref{sec:oracle} supplies the box oracle for that \emph{same} \(K\) using only \textbf{S2}, \textbf{S3}, and effective finite/Presburger constructions. Section~\ref{sec:completion} therefore decides reachability between the original source and target singleton boxes.

Every saturation terminates by the ascending-chain condition for upward sets. Every graph/pattern bound is effective. The construction does not use a presumed algorithm for unrestricted three-counter reachability, a minimum-credit oracle, or semilinearity of the full three-counter relation.

This establishes the stated conditional decidability theorem. It does not establish effective semilinearity of symmetric three-counter reachability; the earlier exponential and quadratic examples in the research log rule out that stronger blanket conclusion.

\section*{Proof record and supporting notes}

The detailed research record is available in the \href{https://chatgpt.com/library?file_id=libfile_f10d80492b688191970c625ab3952096&account_id=personal}{accompanying research log}. The earlier text version is \href{https://chatgpt.com/library?file_id=libfile_7338a158ba3c8191a7b47b1688ce2dd5&account_id=personal}{the assembled proof}.

\paragraph{New main proof ingredients.}

\begin{itemize}[nosep]

\item \path{nov1_grid_pump_completion.md}.

\item \path{nov1_kernel_assembly.md}.

\item \path{nov1_normalization_origin_labels.md}.

\item \path{nov1_profiles.md}.

\item \path{nov1_audit.md} \S\S8--10.

\item \path{nov1_credit.md} \S11.

\item \path{nov1_pair_induction.md} \S\S6--8.

\end{itemize}

\paragraph{Earlier supporting proofs.}

\begin{itemize}[nosep]

\item \path{approach_signed_excursion_normal_form.md}.

\item \path{approach_cone_stratification.md}.

\item \path{positive_block_support_rank.md}.

\item \path{positive_block_integration.md}.

\item \path{finite_s2_s3_axis_pumping.md}.

\item \path{symmetric_2cm_axis_pumping.md}.

\item \path{generalized_axis_perturbation.md}.

\item \path{scalar_oracle_saturation.md}.

\item early research-log high-endpoint/transfer-anchor/diagonal compiler proofs.

\end{itemize}

Historical statements that the global construction was open are superseded only by the explicit normalized-language construction above and its final audits. Failed source-credit, same-face thinness, and raw two-counter fanout arguments are not used.

\end{document}
